\subsection*{Restoring a missing grammatical sense}
One repeated error had a traceable cause in the training data. In the original English gloss corpus, all 211 retained examples of one printed character meant \textit{go}. The source treebank also recorded its frequent use as a third-person object pronoun. A filter requiring the gloss to consist of English letters had discarded the bracketed person annotations. We restored direct-object person supervision and fitted a 51,520-parameter residual before the frozen image speller (Fig.~\ref{fig:pronoun-writer}).

On previously held works the revised writer produced a canonical pronoun on 365/393 marked object occurrences, up from zero, but retained \textit{go} on only four of eight motion occurrences. On held sections of other books it scored 700/723 and 31/40, respectively. A paired test required \textit{go} and \textit{it} for two passages with the same central glyph. The aligned fit completed 4/8 pairs in reserved works and 31/40 in reserved sections. Exchanging answers within the same printed headword during fitting gave 2/8 and 4/40. Thus a preference for the frequent pronoun explains much of the single-answer accuracy but leaves a substantial difference on the paired section test.

In the two independently sourced books, both registered translations were writable at supplied positions on ten of 29 passages, up from five. Four-question performance remained one of 29 with the source reader choosing the positions. In the earlier 539 passages the four-question total rose from 98 to 101. The new supervision changed a grammatical distinction in what the speller could draw. Identifying both requested words, especially when the object precedes the predicate, remains a separate obstacle.
